\section{Use of artificial intelligence}\label{sec:collab}

This paper also serves as a brief case study of the use of AI in mathematical research in actuarial science. This
section describes how AI was used and how the process was handled, and is meant to illustrate a standard for
reporting such use in a mathematical paper. The tools are identified in the declaration before the references. The section also describes how
the models acquire knowledge, the risks this entails for mathematical work, the guardrails used, the formal check
in Lean, and what has and has not been checked.

\subsection{Organization of the work}

The work used Claude (Anthropic) through the Claude desktop application in Cowork mode, which runs the assistant in
a sandboxed environment and can coordinate several AI agents in a workflow. AI agents worked on designated parts of
the project, under the author's direction. A main assistant organized the work and drafted and revised the
manuscript. A roundtable of six agents, five of them each assigned a strand of the literature and one acting as
chair, surveyed the axiomatic theory of premiums, risk measures and performance ratios and produced a synthesis. A
larger panel of agents covered ten literature strands, proposed candidate statements, ranked them by means of judge
agents, drafted proofs, and subjected the proofs to separate AI-assisted checking runs over several rounds. The
literature review was conducted by these AI assistants under the author's supervision: the author defined the
strands to be covered and the questions to be answered, and selected and framed the material presented in
Section~\ref{sec:lit}. Some agents were given the names of researchers as personas, to represent the
perspective of a literature strand; these were simulations, and the researchers concerned did not participate, were
not consulted and bear no responsibility for the content. The research questions, the choice of formulation and all
decisions on the content of the paper were made by the author.

Figures~\ref{fig:workflow} and~\ref{fig:panel} show the organization of the work. In Figure~\ref{fig:workflow}
the author assigns the task of each AI component and reviews its output. Figure~\ref{fig:panel} shows the internal pipeline of the panel, in which candidate statements and
proofs pass through ranking and repeated critical review before being returned to the author.
Figure~\ref{fig:panelpic} gives a schematic picture of the panel: the ten strand agents around a shared digest,
the proof bench of judge, prover and checking agents, and the author's place in the loop.

\begin{figure}[ht]
\centering
\begin{tikzpicture}[font=\footnotesize,
  box/.style={draw,rounded corners=2pt,align=center,inner sep=4pt,minimum height=1.55cm,text width=3.05cm,fill=blue!4},
  hum/.style={draw,thick,rounded corners=2pt,align=center,inner sep=5pt,fill=orange!12},
  arr/.style={-{Stealth[length=2.2mm]},thick},
  sarr/.style={-{Stealth[length=1.8mm]},semithick,gray!70!black}]
\node[box] (rt) {\textbf{Roundtable}\\ 5 literature agents\\ + 1 chair\\ $\to$ synthesis};
\node[box,right=0.45cm of rt] (pn) {\textbf{Panel}\\ 10 literature strands\\ candidate statements\\ judge agents};
\node[box,right=0.45cm of pn] (pr) {\textbf{Proof and review}\\ prover agent\\ AI-assisted checking\\ runs, repair rounds};
\node[box,right=0.45cm of pr] (ma) {\textbf{Main assistant}\\ dialogue, examples,\\ drafting and revision};
\draw[arr] (rt) -- (pn); \draw[arr] (pn) -- (pr); \draw[arr] (pr) -- (ma);
\node[hum,minimum width=14.6cm,above=1.35cm of $(pn.north)!0.5!(pr.north)$] (au)
  {\textbf{Author}: research questions, choice of formulation, decisions on content, review of all outputs};
\foreach \n in {rt,pn,pr,ma}{
  \draw[sarr] ([xshift=-4mm]au.south -| \n.north) -- node[left,font=\scriptsize]{task} ([xshift=-4mm]\n.north);
  \draw[sarr] ([xshift=4mm]\n.north) -- node[right,font=\scriptsize]{output} ([xshift=4mm]au.south -| \n.north);}
\node[hum,below=0.9cm of $(pn.south)!0.5!(pr.south)$,minimum width=6cm] (ms) {\textbf{Manuscript} (author's decisions)};
\draw[arr] (au.east) -- ++(0.35,0) |- (ms.east);
\end{tikzpicture}
\caption{Organization of the work. The author assigns tasks to each AI component and reviews every output; the
manuscript reflects the author's decisions. The stages were used in the order shown, with returns to earlier
stages when the author redirected the work.}\label{fig:workflow}
\end{figure}

\begin{figure}[ht]
\centering
\begin{tikzpicture}[font=\footnotesize,
  st/.style={draw,rounded corners=2pt,align=center,inner sep=3pt,minimum height=1.25cm,text width=2.25cm,fill=blue!4},
  hum/.style={draw,thick,rounded corners=2pt,align=center,inner sep=4pt,text width=2.25cm,minimum height=1.25cm,fill=orange!12},
  arr/.style={-{Stealth[length=2mm]},thick}]
\node[st] (lit) {\textbf{Literature}\\ 10 strand agents};
\node[st,right=0.32cm of lit] (prop) {\textbf{Proposals}\\ candidate\\ statements};
\node[st,right=0.32cm of prop] (jud) {\textbf{Judges}\\ 3 agents rank\\ the candidates};
\node[st,right=0.32cm of jud] (prv) {\textbf{Prover}\\ full statement\\ and proof};
\node[st,right=0.32cm of prv] (chk) {\textbf{Checking}\\ 3 AI-assisted\\ runs};
\node[st,below=1.1cm of chk] (cha) {\textbf{Chair}\\ written note};
\node[hum,below=1.1cm of prv] (aut) {\textbf{Author review}\\ accept, redirect\\ or reject};
\draw[arr] (lit) -- (prop); \draw[arr] (prop) -- (jud); \draw[arr] (jud) -- (prv); \draw[arr] (prv) -- (chk);
\draw[arr] (chk.north) to[bend right=35] node[above,font=\scriptsize]{errors found: repair} (prv.north);
\draw[arr] (chk.south) -- node[right,font=\scriptsize,align=left]{no\\ objections} (cha.north);
\draw[arr] (cha) -- (aut);
\draw[arr,dashed] (aut.west) -- node[above,font=\scriptsize]{new questions and formulations from the author} (aut.west -| lit.south) -- (lit.south);
\end{tikzpicture}
\caption{Pipeline of the panel. Literature agents survey their strands; proposal agents formulate candidate
statements, which judge agents rank. A prover agent drafts a full statement and proof, which separate AI-assisted
checking runs examine; reported errors are sent back for repair over several rounds. A chair agent writes up the
result for the author, who accepts it, redirects the work or rejects it. Agreement among checking runs does not
establish correctness.}\label{fig:panel}
\end{figure}

\begin{figure}[ht]
\centering
\begin{tikzpicture}[font=\scriptsize,
  person/.pic={
    \fill[pic actions] (0,0.30) circle (0.15);
    \fill[pic actions] (-0.24,-0.18) .. controls (-0.24,0.10) and (-0.12,0.13) .. (0,0.13)
                        .. controls (0.12,0.13) and (0.24,0.10) .. (0.24,-0.18) -- cycle;},
  lab/.style={align=center,text width=2.0cm,inner sep=1pt},
  arr/.style={-{Stealth[length=2mm]},thick},
  darr/.style={-{Stealth[length=2mm]},thick,dashed,orange!70!black}]
% table
\fill[brown!12] (0,0) ellipse (2.6 and 1.85);
\draw[brown!60!black,thick] (0,0) ellipse (2.6 and 1.85);
\node[align=center,text width=3.4cm] at (0,0.55) {\textbf{Shared digest}\\[1pt] exact statements with\\ hypotheses, confidence\\ labels and sources\\[1pt] \textbf{+ candidate formulations}};
\pic[fill=teal!70!black,scale=1.15] at (0,-0.95) {person};
\node at (0,-1.45) {\textbf{chair}: synthesis};
% ten strand agents
\foreach \ang [count=\k] in {65.0,90.6,116.1,141.7,167.2,192.8,218.3,243.9,269.4,295.0}{
  \pgfmathsetmacro{\px}{3.05*cos(\ang)} \pgfmathsetmacro{\py}{2.3*sin(\ang)}
  \pic[fill=blue!55!black] at (\px,\py) {person};
  \node[font=\scriptsize\bfseries,blue!55!black] at ($(\px,\py)+(\ang:0.58)$) {\k};
}
% proof bench on the right
\begin{scope}[xshift=6.0cm,yshift=-0.2cm]
  \draw[rounded corners=3pt,fill=gray!8,draw=gray!60] (-1.3,-2.6) rectangle (2.55,2.3);
  \node[font=\footnotesize\bfseries] at (0.6,2.0) {proof bench};
  \foreach \y in {0,1,2} {\pic[fill=purple!60!black] at (-0.45+0.45*\y,1.25) {person};}
  \node[anchor=west] at (0.8,1.3) {judges rank};
  \pic[fill=red!60!black,scale=1.15] at (0,0.05) {person};
  \node[anchor=west] at (0.8,0.05) {prover drafts};
  \foreach \y in {0,1,2} {\pic[fill=green!40!black] at (-0.45+0.45*\y,-1.35) {person};}
  \node[anchor=west,align=left] at (0.8,-1.3) {checking runs\\ try to refute};
  \draw[arr] (0,0.95) -- (0,0.55);
  \draw[arr] (0,-0.35) -- (0,-0.95);
  \draw[arr,red!60!black] (-0.75,-1.3) to[bend left=55] node[left]{repair} (-0.35,0.0);
  \node at (0.6,-2.25) {several rounds};
\end{scope}
\draw[arr] (2.75,0.7) -- node[above,align=center]{candidate\\ statements} (4.65,0.7);
\draw[arr] (4.65,-1.3) -- node[below,align=center]{reviewed\\ result} (2.65,-1.3);
% author
\pic[fill=orange!80!black,scale=1.5] at (2.3,4.9) {person};
\node[align=left,anchor=west] at (2.85,4.9) {\textbf{Author}: chooses the strands and\\ questions, reviews every output\\ and decides the formulation};
\draw[darr] (1.75,4.8) to[bend right=25] node[above left,pos=0.4]{strands, questions} (0.3,3.0);
\draw[darr] (5.6,4.35) to[bend left=15] node[right,pos=0.5,align=left]{accept, redirect\\ or reject} (6.6,2.2);
% key
\node[anchor=north,align=left,text width=14.2cm,font=\scriptsize] at (1.6,-3.15) {%
\begin{tabular}{@{}l@{\hspace{2.5em}}l@{}}
1 insurance-pricing axiomatics & 6 return and star-shaped risk measures\\
2 Choquet integrals and comonotonic additivity & 7 elicitability\\
3 signed Choquet integrals & 8 performance-ratio axiomatics\\
4 law-invariant coherent and convex risk measures & 9 decision theory and functional equations\\
5 acceptability indices and quasi-convex duality & 10 critique of the earlier synthesis
\end{tabular}};
\end{tikzpicture}
\caption{The panel at work. Ten literature agents (numbered as in the key), each assigned a strand by the author, write exact statements
into a shared digest, which together with the candidate formulations forms the knowledge base read in full by the
later agents; a chair agent synthesizes. Candidate statements pass to the proof bench, where judge agents rank them,
a prover agent drafts proofs and separate checking runs try to refute them, with errors returned for repair. The
author sets the strands and questions and accepts, redirects or rejects each result. The icons are schematic: the
agents are separate runs of a language model.}\label{fig:panelpic}
\end{figure}


AI assisted with the following tasks:
\begin{itemize}[nosep]
\item literature exploration;
\item formulation of mathematical statements and drafting of proofs;
\item construction of counterexamples and examples;
\item manuscript drafting and revision.
\end{itemize}


\subsection{How the knowledge base for this paper was assembled}

The agents draw on two kinds of knowledge: what their underlying models learned in training, which can be
recalled only approximately, and documents read into their working context during the task. In this project the
second kind was organized as follows.

The five strands of the roundtable were identified by the author: premium principles and capital allocation;
conic finance and performance ratios; coherent risk measures and acceptance sets; axiomatic risk measures and
elicitability; and the extended-loss-ratio framework with its association. Its synthesis served as the starting
point of the panel.

The panel divided the literature into ten strands: insurance-pricing axiomatics; Choquet integrals and comonotonic
additivity; signed Choquet integrals; law-invariant coherent and convex risk measures; acceptability indices and
quasi-convex duality; return and star-shaped risk measures; elicitability; performance-ratio axiomatics; decision
theory and functional equations; and a critique of the earlier synthesis. Each strand agent returned exact
statements of if-and-only-if theorems with their hypotheses, a confidence label (high, medium or low) and a note
on its source. For example, the signed representation results of
\citet{cmmm2012,cmmm2015} were taken from their working paper versions.

There was no separate retrieval index. Instead, the strand reports were merged into a written digest of about
415{,}000 characters, which, together with the ten candidate formulations (about 245{,}000 characters), formed the
shared knowledge base of the later stages: the judge and prover agents were instructed to read these files in full
before working, and Section~\ref{sec:lit} was written from the digest. The confidence labels therefore
travelled with each statement into the stages that used it.

In addition, the author used Gemini (Google) to prepare summaries of the key works cited in
Section~\ref{sec:lit}: \citet{gerber1974}, \citet{yaari1987} and \citet{wyp1997} on premium principles;
\citet{msm1994}, \citet{cmmm2012,cmmm2015} and \citet{www2020} on signed Choquet integrals; \citet{adeh1999} and
\citet{kusuoka2001} on coherent risk measures; and \citet{chernymadan2009}, \citet{chew1983} and \citet{dekel1986} on
performance ratios and ratio representations. Each summary stated the core results of the work and their role in
this paper. The descriptions in Section~\ref{sec:lit} and the uses of these works in Sections~\ref{sec:axioms},
\ref{sec:results}, \ref{sec:necessity} and~\ref{sec:consequences} were compared with these summaries; no
discrepancy was found, and three descriptions in Section~\ref{sec:lit} were made more precise. Since the summaries
are themselves AI output, this comparison is a consistency check, not a check against the original publications.

Later, the author supplied author versions of three cited papers: \citet{www2020} (working paper of June 2019),
\citet{www2020b} (version of May 2020, with a correction of May 2023) and \citet{hlw2025} (version of June 2024).
The statements taken from them were checked against these versions: the definition of $I_h$, the
characterization theorem used in the representation step and the convexity theorem of \citet{www2020}; the
monotonicity criterion and Lipschitz bound, which are also stated there; the extension to general spaces, the
convention $\mathrm{ES}_1=\operatorname{ess\,sup}$ and the mixture representation of \citet{www2020b}; and the
axioms and the two characterization theorems of \citet{hlw2025}. No discrepancy was found. Lemma numbers that
could differ between these versions and the published articles were removed, and the description of
\citet{hlw2025} in Section~\ref{sec:lit} was rewritten from the paper.

\subsection{Hallucinations and semantic errors encountered}

\emph{Statements about the literature.} The main risk of hallucination concerned the exact hypotheses, conclusions
and numbering of cited results. Two entries of the digest disagreed on whether the mixing measure in the Kusuoka
representation includes an atom at the endpoint; the paper therefore states the mixture over $[0,1]$ with an
explicit convention (Corollary~\ref{cor:shape}(a)) and cites the representation without theorem numbers. For the
same reason, the representation step used in the proofs is not taken from the literature but reproved
(Lemmas~\ref{lem:cashlip}--\ref{lem:rep}).

\emph{Errors in mathematical content.} Table~\ref{tab:errors} lists the errors in AI-generated mathematics that were
found during the project. Most are semantic: the statement proved, or claimed, differed from the statement
intended, typically at a degenerate case (zero premium, the break-even set, the labels $\infty$ and $\star$) or in
the scope of a claim.

\begin{table}[ht]
\centering\footnotesize
\begin{tabularx}{\textwidth}{@{}>{\raggedright\arraybackslash}X>{\raggedright\arraybackslash}p{2.6cm}>{\raggedright\arraybackslash}X@{}}
\toprule
\textbf{Error} & \textbf{Found by} & \textbf{Resolution}\\
\midrule
$\E/C_g$ with concave $g$ claimed quasi-convex (early synthesis) & panel agents & quasi-concave; Cor.~\ref{cor:shape}(d)\\
quasi-convexity stated as an if-and-only-if in four candidate formulations & judge agent ($h=t$, $g=\sqrt t$, $\beta=2$) & only sufficient conditions kept\\
candidate route via Sincov and Cauchy equations & judge agent (continuous $h$ not of bounded variation) & route abandoned\\
main argument claimed to apply verbatim when $(\rho,\Pi)$ is only assumed to exist & checking run & separate argument given\\
association claimed to hold exactly on the finite domain & checking run & holds off $\{\rho=\Pi=0\}$; Cor.~\ref{cor:assoc}(a)\\
``$h-g$ concave'' as a sufficient condition for quasi-convexity & main assistant, during rewriting & $h$ concave, $g$ convex\\
converse of Theorem~\ref{thm:A} failed for \axLI{} (labels unconstrained) & checking run & \axLI{} compares finite values only\\
claimed equivalence with the cash-line formulation & checking run & convention \ax{S} added; Remark~\ref{rem:oldCP}\\
regularity axioms undefined for two counterexamples & checking run & \axRC, \axRM{} stated for $\Pi$ and $\rho_D$\\
uniqueness of the extension in Prop.~\ref{prop:bridge} & review with ChatGPT & unique up to labels, unique under \ax{S}\\
independence claimed for (A0) and (A1) & review with ChatGPT & treated as illustrations; Sec.~\ref{sec:necessity}\\
mixture of expected shortfalls missing the endpoint & review with ChatGPT & $\mathrm{ES}_1:=\operatorname{ess\,sup}$; Cor.~\ref{cor:shape}(a)\\
\bottomrule
\end{tabularx}
\caption{Errors in AI-generated mathematics found during the project, and their resolution.}\label{tab:errors}
\end{table}

\subsection{Guardrails applied in this project}

\begin{itemize}[nosep]
\item \emph{Labels on knowledge.} Every literature statement in the digest carried a confidence label and a note on
  its source, as described above.
\item \emph{Independent summaries.} Summaries of the key cited works prepared with Gemini were compared with their
  descriptions and uses in the paper.
\item \emph{Self-contained proofs.} Results whose exact statements were uncertain are reproved; in particular the
  representation of comonotonic additive functionals, for which the proof cites \citet[Thm.~1]{www2020}, is
  reproved in Lemmas~\ref{lem:cashlip}--\ref{lem:rep}.
\item \emph{Adversarial checking.} The three checking runs on the panel's proofs were instructed to refute; one tested
  the degenerate case $\LR\equiv1$ with a premium that is not comonotonic additive, and two recomputed the numbers in
  the counterexamples. The restructured formulation of Sections~\ref{sec:axioms}--\ref{sec:proofs} was given to a
  further checking run with instructions to look for circularity and for the cases $\Pi(X)=0$, break-even books,
  $\LR\equiv1$ and $c=0$.
\item \emph{Repair loop.} Reported gaps were returned to the prover agent; the panel's proofs went through two
  repair rounds before the final checking round.
\item \emph{Counterexamples.} Each axiom other than \axP{} is tested against an explicit loss ratio
  (Section~\ref{sec:necessity}).
\item \emph{Recomputation.} The values in Examples~\ref{ex:PH}, \ref{ex:ES}, \ref{ex:qc} and in the counterexamples
  were recomputed by script, and the monotonicity range in Example~\ref{ex:family} was confirmed on a grid.
\item \emph{Formal checking.} The loss-ratio-specific steps were formalized in Lean and compiled; see
  Section~\ref{sec:collab-lean}.
\item \emph{A second system.} The review with ChatGPT, a model from a different provider, found three of the errors in
  Table~\ref{tab:errors} that the earlier checks had missed.
\end{itemize}

\subsection{What the prover produced and what was formally checked}

The ``prover'' of Figure~\ref{fig:panel} is a language-model agent that writes proofs in ordinary mathematical
language; it is not a formal theorem prover. In this project it drafted the first complete if-and-only-if theorem,
in a formulation based on an assumed cash-line axiom, with proofs of the results that are now
Theorems~\ref{thm:A}, \ref{thm:B}, \ref{thm:A0}, \ref{thm:C} and~\ref{thm:Cp} and of the appendix lemmas; the final
checking round of these proofs returned two reports without objections and one report of gaps in side claims, which
were repaired. The present formulation through the implied premium and risk measure
(Lemmas~\ref{lem:premium}--\ref{lem:LI} and Proposition~\ref{prop:readings}) was drafted afterwards by the main
assistant and checked by one further run, whose findings appear in Table~\ref{tab:errors}.

The loss-ratio-specific steps were also checked formally in Lean (Section~\ref{sec:collab-lean}).

\subsection{Formal checking in Lean}\label{sec:collab-lean}

\subsubsection*{Lean and what a formal check establishes}

Lean~4 \citep{demoura2021} is an interactive proof assistant: a programming language in which definitions,
statements and proofs are written in a formal language based on dependent type theory, and in which every proof is
checked by a small trusted program, the kernel. A proof may be written step by step or with the help of tactics,
procedures that search for parts of a proof automatically, but whatever a tactic produces is checked again by the
kernel. Mathlib \citep{mathlib2020} is the community library of formalized mathematics for Lean; it supplies the
real numbers, groups, limits and the other notions used below.

Two features of Lean are used to state precisely what has been checked. A proof can be left unfinished by writing
\texttt{sorry}, which Lean accepts with a warning; a finished development contains none. The command
\texttt{\#print axioms} lists the foundational axioms on which a theorem depends; ordinary classical mathematics in
Lean uses exactly three of them, \texttt{propext} (propositional extensionality), \texttt{Classical.choice} (the
axiom of choice) and \texttt{Quot.sound} (soundness of quotients), and an unfinished proof shows up in this list as
\texttt{sorryAx}.

A successful check establishes that the formal statements follow from the formal definitions. It does not
establish that the formal definitions and statements say what the paper says, which must be checked by a human
reader, and it says nothing about results that enter as hypotheses or are not formalized.

For this paper a formal check is useful for three reasons. First, the proofs were drafted with AI assistance, and
the errors recorded in Table~\ref{tab:errors} show that such drafts can contain steps that read well and are wrong;
the kernel does not depend on how a proof was produced. Second, the steps that are specific to loss ratios
(Lemmas~\ref{lem:premium}--\ref{lem:LI}) are elementary but delicate, with case distinctions at $\Pi(X)=0$, at
break-even books and at the price point, which is where informal checking is most likely to miss a case. Third,
formalization forces the hypotheses to be stated exactly, and here it shows that those steps use only a few
algebraic properties of comonotonicity and equality in law (see ``What was formalized'' below).

\subsubsection*{What was formalized}\label{sec:lean-scope}

The formalization works in an abstract setting. Positions form an additive commutative group $\Xp$ with an additive
embedding $c\mapsto\mathrm{const}(c)$ of the sure losses, a relation ``comonotone'' under which every position is
comonotone with every constant and which is preserved when both positions are shifted by constants, and a relation
``equal in law'' preserved by a common cash shift. A loss ratio is a map from $\Xp$ to $\R\cup\{\mathtt{none}\}$,
where \texttt{none} stands for a non-finite value, so that the labels $\infty$ and $\star$ are identified as in
\axLI. The axioms read as follows; \texttt{prem} is the implied premium and \texttt{rhoD} is $\Pi\cdot\LR$.

{\small
\begin{verbatim}
def AxP : Prop := ∀ x : X, ∃! c : ℝ, LR (x + S.const c) = none
def AxN : Prop := ∀ c : ℝ, c ≠ 0 → LR (S.const c) = some 1
def AxLI : Prop := ∀ x y : X, S.eqd x y → LR x = LR y
def AxCM (hP : AxP S LR) : Prop :=
  ∀ x y : X, S.comon x y → LR x ≠ none → LR y ≠ none →
    prem hP x + prem hP y ≠ 0 →
    LR (x + y) = some ((rhoD hP x + rhoD hP y) / (prem hP x + prem hP y))
\end{verbatim}
}

These coincide with Axioms~\ref{ax:P}, \ref{ax:N}, \ref{ax:CM} and~\ref{ax:LI}. The file has about 500 lines with
10 definitions and 26 lemmas and theorems. Table~\ref{tab:lean} lists the results of the paper that are formalized.

\begin{table}[ht]
\centering\small
\begin{tabular}{@{}ll@{}}
\toprule
\textbf{Result in the paper} & \textbf{Lean name}\\
\midrule
Implied premium, uniqueness and cash additivity (Lemma~\ref{lem:premium}) & \texttt{prem}, \texttt{prem\_shift}, \texttt{prem\_const}\\
Implied risk measure and cash additivity (Lemma~\ref{lem:implied}) & \texttt{rho}, \texttt{rho\_cash}\\
Möbius form and residue (Proposition~\ref{prop:readings}(a)--(c)) & \texttt{mobius}, \texttt{rho\_eq}, \texttt{residue}\\
Comonotonic additivity, also on the break-even set (Lemma~\ref{lem:comon}) & \texttt{comon\_additive}\\
Law invariance of the implied pair (Lemma~\ref{lem:LI}) & \texttt{LI\_transfer}\\
Converse construction and uniqueness of the pair & \texttt{converse}, \texttt{unique\_pair}\\
Association (Corollary~\ref{cor:assoc}(a)) & \texttt{association}\\
Pooling at non-finite values (Corollary~\ref{cor:projpool}) & \texttt{pooling\_extended}\\
Algebraic core of Proposition~\ref{prop:F2} & \texttt{no\_monotone\_core}\\
Theorem~\ref{thm:A}, given the representation step & \texttt{main\_conditional}\\
\bottomrule
\end{tabular}
\caption{Results formalized in Lean.}\label{tab:lean}
\end{table}

The theorem \texttt{main\_conditional} takes the representation step as a hypothesis: any cash-additive,
comonotonic additive, law-invariant and regular functional is of the required form \citep[Thm.~1]{www2020}. From
\axP, \axN, \axCM{} and \axLI{} it derives that the implied premium and the implied risk measure satisfy the
hypotheses of that step and that $\LR=\bracket{\rho}{\Pi}$ everywhere. Not formalized are the representation step
itself (Lemmas~\ref{lem:cashlip}--\ref{lem:rep}), which would require a theory of Choquet integrals that Mathlib does not
provide, the fact that $L^\infty$ with comonotonicity and equality in law is an instance of
the abstract setting, the layer formula, and the examples and counterexamples.

\subsubsection*{How the check was carried out}

The Lean file was written by the main assistant. Lean could not be installed in the assistant's cloud workspace,
whose network policy blocks the servers that distribute Lean, so the file was compiled on the author's computer.
The author installed Lean and created the project; the assistant read the build output from the project folder
through the link between the Claude desktop application and the author's computer, wrote corrected versions of the
file into the same folder, and the author ran the next build (Figure~\ref{fig:leanloop}).

\begin{figure}[ht]
\centering
\begin{tikzpicture}[node distance=6mm and 14mm,
  box/.style={draw,rounded corners,align=center,font=\small,minimum height=10mm,text width=34mm},
  >={Stealth[length=2mm]}]
\node[box] (a) {Assistant writes or\\corrects \texttt{Basic.lean}};
\node[box,right=of a] (b) {File written into the\\author's project folder};
\node[box,below=of b] (c) {Author runs\\\texttt{lake build}};
\node[box,left=of c] (d) {Assistant reads\\\texttt{build\_log.txt}};
\draw[->] (a) -- (b);
\draw[->] (b) -- (c);
\draw[->] (c) -- (d);
\draw[->] (d) -- node[left,font=\footnotesize,align=center]{errors}(a);
\node[box,right=of c,text width=30mm] (e) {Build succeeds;\\axiom report};
\draw[->] (c) -- (e);
\end{tikzpicture}
\caption{The build-and-repair loop used for the Lean check.}\label{fig:leanloop}
\end{figure}

Table~\ref{tab:leantime} gives the timeline. Times are in UTC and are taken from the session records and from the
modification times of the files in the project folder.

\begin{table}[ht]
\centering\small
\begin{tabular}{@{}>{\raggedright\arraybackslash}p{3.3cm}>{\raggedright\arraybackslash}p{11.0cm}@{}}
\toprule
\textbf{Time (UTC)} & \textbf{Step}\\
\midrule
21 Oct, 21:38 & The assistant writes a first Lean file for the results in Table~\ref{tab:lean}; it cannot be compiled
  in the cloud workspace.\\
24 Oct, 17:08 & The author asks how a formal proof check can be used; the assistant describes what is feasible
  (the loss-ratio steps), what is hard (the representation step) and how to install Lean.\\
24 Oct, 17:26--17:31 & The author installs the Lean extension for the code editor and the Lean toolchain manager.\\
24 Oct, 17:35--17:37 & The author creates the project \texttt{loss\_ratio\_lean} with Mathlib; Lean and Mathlib
  version~4.36.0-rc1; the precompiled Mathlib files are downloaded.\\
24 Oct, 17:38--17:39 & The assistant inspects the project configuration, confirms that it is complete, and places
  the Lean file in it, replacing the template placeholder.\\
25 Oct, 12:47--12:53 & First attempt to build: the command \texttt{lake} is not found, because the toolchain had not
  been added to the terminal's search path; fixed by loading the toolchain environment.\\
25 Oct, 12:55 & Build 1: four errors, none in a definition or statement. One proof stopped one algebraic step short;
  in three proofs Lean's simplifier rewrote $\mathrm{const}(-c)$ as $-\mathrm{const}(c)$, so that cash additivity
  no longer applied.\\
25 Oct, 12:57 & The assistant writes a corrected file and appends \texttt{\#print axioms} for ten main theorems.\\
25 Oct, 13:04 & Build 2: one error, in the converse construction, where the goal was displayed in an unexpanded
  form. The axiom report shows the three standard axioms for nine theorems and \texttt{sorryAx} for the converse.\\
25 Oct, 13:05 & The assistant adds one line that states the goal explicitly.\\
25 Oct, 13:08 & Build 3 succeeds. All ten theorems depend only on the three standard axioms.\\
25 Oct, 13:09 & The assistant compares the Lean definitions of \axP, \axN, \axCM{} and \axLI{} and the theorem
  statements with Section~\ref{sec:axioms}.\\
25 Oct, 13:37 & At the author's request, the project's own build output is deleted and the file is rebuilt from
  scratch: success in 8.8\,s, with the same axiom report.\\
\bottomrule
\end{tabular}
\caption{Timeline of the Lean check.}\label{tab:leantime}
\end{table}

\subsubsection*{Checks performed}

The following checks were made to ensure that the outcome is what is reported.

\begin{itemize}[nosep]
\item \emph{No change to statements during repair.} The corrections after builds 1 and 2 changed proof steps only;
  no definition and no theorem statement was altered.
\item \emph{No unfinished proofs.} The final file contains no \texttt{sorry}, no \texttt{admit} and no added
  \texttt{axiom}.
\item \emph{Axiom report.} \texttt{\#print axioms} for each of the ten theorems in Table~\ref{tab:lean} lists only
  \texttt{propext}, \texttt{Classical.choice} and \texttt{Quot.sound}.
\item \emph{Clean rebuild.} The final build started from an empty build directory, so its result does not rely on
  cached output.
\item \emph{Identity of the file.} The file compiled in the final build was compared with the file reported here and
  found identical.
\item \emph{Fidelity of statements.} The formal axioms and theorem statements were compared with the paper by the
  assistant; this comparison is to be repeated by the author before the code is released.
\end{itemize}

The remaining warnings concern style and the use of renamed library lemmas; they do not affect correctness. The Lean
source and the build log are available in the repository
\url{https://github.com/hirbod-assa/loss-ratio-lean}, which contains the Lean project (Lean and Mathlib
version~4.36.0-rc1) together with instructions for rebuilding it with \texttt{lake build}.


\subsection{Verification status}\label{sec:collab-status}

The proofs have undergone AI-assisted critical review. Such review is fallible and does not
constitute formal verification; this applies equally to the informal review with ChatGPT. No claim of
machine-checked correctness is made beyond the results formalized in Lean, in the abstract setting and with the
hypothesis described in Section~\ref{sec:collab-lean}. Independent human review of all proofs, and journal peer review, have not yet taken
place, and apart from the three papers named above the references have not yet been checked against the original
publications.
